\section{Tổng kết}
\label{sec:8-1}

Dự án đã biến 568.454 review thực phẩm Amazon thành một bộ dữ liệu 393.577 review sạch, không rò rỉ (Chương~\ref{ch:data-preparation}), chia 80/10/10 có phân tầng, và so sánh hai bộ phân loại đã huấn luyện với quy tắc luôn đoán positive trong bài toán cảm xúc ba lớp gán nhãn từ số sao.

\begin{itemize}
    \item \textbf{BERT-base được fine-tune là mô hình tốt hơn:} Macro-F1 0,7877 và accuracy 0,9074 trên 39.358 review test, hơn BiLSTM với attention (0,7422) là 0,0455, với khoảng 95\% của chênh lệch là $[+0{,}0400;\,+0{,}0507]$. Nó hơn ở cả ba lớp, nên học chuyển giao tốt hơn học từ đầu trong bài toán này.
    \item \textbf{Neutral là điểm yếu của cả hai mô hình} (F1 0,49 và 0,55): chúng dự đoán nó quá nhiều, và dịch ngưỡng chỉ giúp BiLSTM.
    \item \textbf{Chi phí.} BERT cần gấp 19 lần thời gian huấn luyện và 17,5 lần số tham số của BiLSTM để có thêm 0,0455 Macro-F1.
\end{itemize}

\section{Hạn chế}
\label{sec:8-2}

Các hạn chế được liệt kê cùng số đo ở Mục~\ref{sec:7-9}. Quan trọng nhất là: nhãn sao không phải lúc nào cũng khớp với văn bản; cách chia ngẫu nhiên không đo hiệu năng trên review tương lai; mỗi mô hình chỉ chạy một lần; và mô hình BERT chỉ được fine-tune hai epoch, với Macro-F1 validation vẫn đang tăng.

\section{Hướng phát triển}
\label{sec:8-3}

\begin{enumerate}
    \item \textbf{Fine-tune BERT lâu hơn hoặc lặp lại} (nhiều epoch hơn, nhiều seed), và các bộ mã hoá mạnh hơn như RoBERTa-base hay DeBERTa-v3, để kiểm tra điểm của BERT có phải là cận dưới không.
    \item \textbf{Chia theo thời gian} (huấn luyện đến 2011, kiểm tra trên 2012) để đo hiệu năng trên review sau này và ảnh hưởng của việc phân bố nhãn trôi.
    \item \textbf{Vector từ tiền huấn luyện} (GloVe, fastText) cho embedding của BiLSTM, để xem khoảng cách với BERT có thu hẹp không.
    \item \textbf{Các cách xử lý mất cân bằng khác} (focal loss, oversampling), vì trọng số neutral lớn có thể gây việc dự đoán neutral quá nhiều.
    \item \textbf{Cắt đầu và đuôi} cho review dài, nơi cả hai mô hình đạt điểm thấp nhất.
    \item \textbf{Một mẫu nhỏ do con người gán nhãn}, không dựa vào số sao, để đo bao nhiêu lỗi là do nhãn nhiễu.
\end{enumerate}

\section{Khả năng tái lập}
\label{sec:8-4}

Cách chia dùng seed và quy tắc cố định, và mọi dự đoán trên tập test đều được lưu, nên mọi bảng ở Chương~\ref{ch:evaluation} có thể tính lại từ đó. Mã nguồn gồm một notebook (làm sạch, tìm hiểu và chuẩn bị dữ liệu, BiLSTM, BERT và đánh giá) và một gói dùng chung \texttt{src/}.
